% WACV 2027 author response / rebuttal for the MMCTL manuscript.
% Point-by-point response: each comment is quoted, followed by a Response that
% lists the specific manuscript additions. Compile with pdflatex + bibtex.
\documentclass[10pt,twocolumn,letterpaper]{article}
\usepackage[rebuttal,algorithms]{wacv}  % Algorithms track (matches the main paper)

\usepackage{graphicx}
\usepackage{amsmath}
\usepackage{amssymb}
\usepackage{booktabs}
\usepackage{enumitem}

%
% --- inline annotations
%
\newcommand{\red}[1]{{\color{red}#1}}
\newcommand{\todo}[1]{{\color{red}#1}}
\newcommand{\TODO}[1]{\textbf{\color{red}[TODO: #1]}}
% --- disable by uncommenting  
% \renewcommand{\TODO}[1]{}
% \renewcommand{\todo}[1]{#1}




\definecolor{wacvblue}{rgb}{0.21,0.49,0.74}
\usepackage[pagebackref,breaklinks,colorlinks,allcolors=wacvblue]{hyperref}

\def\wacvPaperID{*****}
\def\confName{WACV}
\def\confYear{2027}

% Compact list + response macros so the letter reads uniformly.
\setlist[itemize]{nosep,leftmargin=1.35em,topsep=2pt}
\newcommand{\cmt}[2]{\smallskip\noindent\textbf{Comment #1.}\ \textit{#2}\par}
\newcommand{\resp}[1]{\noindent\textbf{Response.}\ #1\par}

\begin{document}

\title{Response to Reviewers:\\ Learning Annotation-Driven Latent Structure for Multi-Modal Fusion}

\maketitle
\thispagestyle{empty}

\noindent We thank both reviewers for their careful reading and constructive
suggestions. We have revised the manuscript accordingly; all newly added content
is marked in \red{red} in the revised PDF for ease of review. Below we respond to
each comment and indicate the specific additions. Equation, figure, and table
numbers refer to the revised manuscript.

%========================================================================
\section{Reviewer 1}

\cmt{1}{Clarify the essential difference between MMCTL and supervised contrastive
learning or cross-modal triplet loss, especially the added contribution of
continuous similarity constraints.}
\resp{We rewrote the Related Work section and added an explicit statement of the
contribution. Unlike binary supervised contrastive learning and fixed-margin
cross-modal triplets, MMCTL orders embedding distances by graded, continuous
annotation proximity rather than binary co-membership. Concretely, MMCTL adds:}
\begin{itemize}
  \item a joint annotation vector with feature-wise normalization and per-task weights (Eqs.~1--3);
  \item a margin-free continuous triplet whose separation target is the annotation-distance gap (Eq.~7), with no tuned margin;
  \item a pairwise distance-alignment regularizer matching embedding to annotation distances (Eq.~8);
  \item extension to a multi-modal shared space where same-object, different-modality pairs have zero annotation distance, directly pulling cross-modal embeddings together.
\end{itemize}

\cmt{2}{The paper claims MMCTL learns a more interpretable latent space;
visualizations or cross-modal retrieval examples are recommended.}
\resp{We added embedding-space visualizations and quantitative interpretability metrics:}
\begin{itemize}
  \item Fig.~1: AWIR fused embeddings (PCA and t-SNE) across all fusion methods, colored by class and sized by bounding-box area;
  \item Fig.~3: per-class ``zoom'' panels colored by the continuous label, showing smooth within-class transitions;
  \item Table~1: standard metrics---silhouette (class separation) and Spearman distance-preservation (continuous attributes).
\end{itemize}

\cmt{3}{Specify the rules for selecting positive and negative samples in
continuous triplets, as well as the similarity thresholds or margin settings.}
\resp{Clarified in the Shared Latent Space subsection. We use a \emph{batch-all}
strategy: every in-batch index is an anchor, and a pair $(p,q)$ is a valid
(positive, negative) pair iff $p$ is closer to the anchor than $q$ in joint
annotation space. Positives and negatives are therefore defined solely by
annotation-distance ordering; there is \emph{no} hard similarity threshold and
\emph{no} tuned margin (the loss is margin-free).}

\cmt{4}{How are annotation attributes such as category, object size, and height
normalized and weighted? Does the weighting strategy affect the final results?}
\resp{The Problem Setup gives feature-wise min--max normalization, task-length
rescaling, and fixed per-task weights (Eqs.~1--3; AWIR $w_{\text{class}}{=}0.5$,
NEON $w_{\text{class}}{=}1.0$). We added a class-weight sensitivity study:}
\begin{itemize}
  \item Table~3 and Fig.~4 sweep $w_{\text{class}}\in\{0.5,0.8,1.0,1.5\}$;
  \item accuracy rises monotonically ($74.9\%\rightarrow84.6\%$) while box/height regression stays stable, confirming the weight acts as an interpretable semantic-vs-geometric trade-off.
\end{itemize}

\cmt{5}{What is the rationale for freezing the encoder in the two-stage training
scheme? A comparison with end-to-end fine-tuning would help validate this.}
\resp{We justify freezing as a fair fusion-evaluation protocol (it isolates
fusion geometry from unequal encoder capacity) and added a controlled comparison
on NEON from the same pretrained encoders:}
\begin{itemize}
  \item frozen $81.99\pm1.02\%$ vs.\ end-to-end $80.66\pm1.35\%$ accuracy;
  \item end-to-end yields only marginal gains on the continuous heads, so freezing is justified and substantially cheaper.
\end{itemize}

\cmt{6}{The baseline set is limited. Representative multimodal methods such as
supervised contrastive learning, cross-attention, or Transformer-based fusion
should be included.}
\resp{We removed the DeepSet baselines and added three stronger, representative
fusion baselines with citations: supervised contrastive fusion (SupCon),
cross-attention pooling, and a Transformer encoder over modality tokens. They
appear in Tables~1 and~2 and in the embedding visualizations.}

\cmt{7}{Report standard deviations and statistical significance tests to verify
whether MMCTL's improvements are stable and reliable.}
\resp{All results tables now report mean\,$\pm$\,std over the cross-validation
folds, and we added paired one-sided Wilcoxon signed-rank tests:}
\begin{itemize}
  \item MMCTL's continuous-geometry advantage is significant on AWIR position ($p{=}0.0078$), with a positive trend on box ($p{=}0.098$);
  \item we are explicit that MMCTL's contribution is balanced coherence across metrics rather than winning every individual metric.
\end{itemize}

%========================================================================
\section{Reviewer 2}

\cmt{1}{State where the substantive increment relative to prior work lies---the
joint annotation vector/weighting (Eqs.~1--3), the margin-free continuous triplet
(Eq.~7), the distance-alignment regularizer (Eq.~8), or the extension from
single-modal to multi-modal shared latent space.}
\resp{The primary increment is the extension to a \emph{multi-modal} shared
latent space in which same-object, different-modality pairs share an identical
joint annotation vector (zero annotation distance), directly coupling cross-modal
embeddings. This is accompanied by the joint annotation vector
construction/weighting (Eqs.~1--3), the margin-free continuous triplet (Eq.~7),
and the distance-alignment regularizer (Eq.~8). Related Work now states this
explicitly relative to our prior unimodal work.}

\cmt{2}{The introduction uses gradient conflict and task dominance as motivation
but provides no supporting experiments; revise the introduction to highlight the
problem the paper solves.}
\resp{We revised the Introduction to remove the unsupported gradient-conflict /
task-dominance motivation and instead frame the problem directly: task-driven
fusion shapes latent geometry only indirectly, so embedding distances are not
constrained to reflect annotation-defined similarity. The motivation now matches
the problem addressed and the evidence provided.}

\cmt{3}{Discuss the differences between MMCTL and existing soft/continuous
similarity metric learning and adaptive-margin triplet losses.}
\resp{Related Work now contrasts MMCTL with continuous/soft-similarity metric
learning (Kim \emph{et al.}) and adaptive-margin triplets (Ha \emph{et al.}):
MMCTL is margin-free (its separation target is the annotation-distance gap, with
no tuned margin) and operates in a multi-modal shared space, which those methods
do not.}

\cmt{4}{MMCTL also uses task labels; the difference from baselines is the loss
form, not the information used. Add cross-dataset/cross-region transfer or
missing-modality experiments.}
\resp{We agree the distinction is the inductive bias (loss form), not the label
information, and we now state this explicitly. We added missing-modality
experiments:}
\begin{itemize}
  \item zeroing one modality half of the MMCTL joint embedding and re-probing yields graceful degradation (drop-thermal $84.7\%$, drop-RGB $86.2\%$) that stays above the unimodal baselines ($\sim$78--80\%), evidencing complementary but non-redundant modality use;
  \item cross-region transfer is noted as future work in the Conclusion.
\end{itemize}

\cmt{5}{The batch-all strategy is used but the number of triplets and
computational cost are not reported. Report total triplets, training time per
epoch, and GPU memory usage.}
\resp{We added a triplet-compute analysis. With batch size $32$ and modality
concatenation ($64$ embeddings), the batch-all construction yields at most
$249{,}984$ ordered triplets per batch ($\sim$2.0M per epoch) before the
annotation-order mask; a microbenchmark measures $\sim$1.7\,s per training step
(batch $8$) on an Apple M4 CPU. Our runs were CPU-based, so peak GPU memory was
not measured; we state this explicitly and will add GPU timing/memory if GPU logs
are required.}

\cmt{6}{Clarify the relationship between ``k-fold cross-validation within the
training set'' and ``averaged over 5 runs'' in the caption of Figure~3.}
\resp{Clarified in the Experimental Setup: the ``runs'' reported in the captions
\emph{are} the cross-validation folds ($8$ for AWIR, $5$ for NEON), and each
mean\,$\pm$\,std is computed across these folds. The ambiguous ``runs'' wording
has been reconciled with the k-fold description.}

\cmt{7}{Some downstream results need more detailed discussion: add error analysis
of typical failure cases and explain when the method is effective and when it is
not.}
\resp{We expanded the downstream discussion to explain when MMCTL helps and when
it does not:}
\begin{itemize}
  \item class-supervised baselines win raw class separation, and a directly sensing modality (LiDAR) wins its own continuous attribute, but each fails the other axis; MMCTL is the only method balanced across both;
  \item the per-class within-class analysis (Fig.~3) identifies classes with strong continuous transitions (e.g., PIPA2 canopy height) versus weaker ones, providing a class-level effectiveness/error analysis.
\end{itemize}

\cmt{8}{Update references to their formally published versions, e.g.,
reference~[8].}
\resp{We are updating the bibliography to cite formally published versions in
place of arXiv preprints where available, including reference~[8]. This audit is
being finalized against the compiled reference list.}

\end{document}
